At a point of $E$ with the collar constants $\mu, \nu$ and the bump held fixed, call a vector $u'$ with $k(u') = 1$ a \emph{candidate value} for the dying direction $k$; the candidates form an affine space over $\ker k$. For every $\lambda \in \ker k$ the composite $\theta_\lambda := \hat g_\lambda \circ \Phi_\lambda^{-1}$, with $\hat g_\lambda := g_\lambda \oplus \id$ on $T = \shW \oplus \operatorname{span}\{\nO, \nInf\}$, is an automorphism of the completed structure $(T, h_u, k)$: an $h_u$-isometry restricting to $g_\lambda$ on $\shW$ and fixing $k$, hence an automorphism of the completed Clifford algebra preserving its grading and its degeneration subalgebra. Its contragredient action on vectors, $(\theta_\lambda|_{\shW})^{-\mathsf{T}} = \id - \lambda \otimes k$, translates every candidate by $-\lambda$. A natural assignment of a candidate value, one per completed structure and commuting with automorphisms, would therefore be a fixed point of a free translation action of $(\ker k, +)$: none exists. Nor is a ladder partner for $k$ any more canonical than a value: for $\lambda \neq 0$, $h_u(k, x) = 0$ for every $x \in \operatorname{Fix}(\theta_\lambda)$, whereas a ladder partner has $h_u(k, x) = \tfrac12$; and the construction attaches to each candidate $u' = u + v$ the partner $\Phi_v(k^*) = k^* + \tfrac12\tilde v$ of the transversal-independence theorem, injectively in $v$. An analogue of the no-repair theorem of \cite[Thm.~4.2]{ButlerT1} thus holds at the wall itself, where the base form carries no hyperbolic partner for $k$ and that theorem's hypotheses fail.
\emph{Provenance:} \cite[Thm.~10.4]{ButlerT4} (\claimref{sheaf.no-equivariant-value-wall}), verbatim up to the recorded notational alignment ($E$ for that paper's $\Sigma$, $\shW$ for its $T^*M$). The objects in the statement are that paper's: the \emph{completed structure} $(T, h_u, k)$ is the ambient module $T$ with the ambient form $h_u$ of the coupling construction built on the transversal $u$ (Fact \claimref{singularity.imported-one-sided-completion-singularity}) and the dying direction $k$, carrying the $\mathbb{Z}_2$-grading of its Clifford algebra and the subalgebra generated by $\shW$; the named frame $(\nO, \nInf)$ and the partner $k^* = \nO/2\nu$ are choices inside it, not part of it. Its supporting results are \cite[Lems 10.1, 10.2, Thm.~10.3]{ButlerT4} (\claimref{sheaf.base-unipotents-wall}, \claimref{sheaf.transversal-independence-wall}, \claimref{sheaf.solve-closed-radical-gauge}): the base unipotents $g_\lambda(\alpha) = \alpha + \alpha(\lambda)k$, $G$-isometries fixing $k$; for $v \in \ker k$ the representative $\tilde v \in \shW$ with $G(\alpha, \tilde v) = \alpha(v)$, normalised by $G(\tilde v, \tilde v) + 2\tilde v(u) = 0$; and the unipotent intertwiner $\Phi_v$, the identity on $\shW$ with $\Phi_v(\nO) = \nO + \nu\tilde v$, $\Phi_v(\nInf) = \nInf + \mu\tilde v$, an isometry $(T, h_{u+v}) \to (T, h_u)$ satisfying the cocycle $\Phi^u_v \circ \Phi^{u+v}_w = \Phi^u_{v+w}$, so that the completions built on all transversals are canonically identified, at the wall, by a coherent system of isomorphisms preserving the form, the grading and the degeneration subalgebra but neither the frame nor the ladder pair. Off the wall the two ambient forms need not be isometric (the shadow wall moves with the transversal); nothing off the wall is used. All of it is machine-verified there on the Schwarzschild boundary data at two masses.
